\documentclass[conference]{IEEEtran}
\IEEEoverridecommandlockouts
% Remove \IEEEoverridecommandlockouts if no funding footnote is needed.

\usepackage{cite}
\usepackage{amsmath,amssymb,amsfonts}
\usepackage{graphicx}
\usepackage{textcomp}
\usepackage{xcolor}
\usepackage{booktabs}
\usepackage{multirow}
\usepackage{url}

\def\BibTeX{{\rm B\kern-.05em{\sc i\kern-.025em b}\kern-.08em
    T\kern-.1667em\lower.7ex\hbox{E}\kern-.125emX}}

\begin{document}

\title{Consensus-Driven Evaluation of Medical AI Answers:
A Privacy-Preserving Multi-SLM Judge Framework}

\author{
\IEEEauthorblockN{First Author}
\IEEEauthorblockA{\textit{Department Name} \\
\textit{University / Institution}\\
City, Country \\
email@domain}
\and
\IEEEauthorblockN{Second Author}
\IEEEauthorblockA{\textit{Department Name} \\
\textit{University / Institution}\\
City, Country \\
email@domain}
}

\maketitle

\begin{abstract}
Automated evaluation of AI-generated medical answers typically relies on
a single large cloud-hosted model, posing HIPAA compliance risks and
precluding offline deployment. We present Multi-SLMs-as-Judge, a
privacy-preserving framework that assembles a panel of four locally
deployed small language models (SLMs) to evaluate medical answer quality.
Inter-judge consensus is quantified via a normalized pairwise agreement
score, which gates quality score reporting and routes low-consensus cases
to human review. Evaluated on 100 clinical questions spanning five domains
and five rubrics, binary rubrics achieve mean agreement of 91--94\%, while
Likert-scale rubrics fall to 72--81\%. A controlled binary-vs-Likert
comparison on identical criteria isolates a 22 percentage-point drop
attributable solely to scale choice. Cardiology is the most consistently
judged domain; pharmacology generates the most disagreement, driven
primarily by Likert rubrics. These findings give practical guidance for
rubric design and scoring scale selection in privacy-safe clinical
evaluation pipelines.
\end{abstract}

\begin{IEEEkeywords}
medical AI evaluation, small language models, inter-judge agreement,
privacy-preserving NLP, rubric design, local deployment, e-health
\end{IEEEkeywords}

% ===================================================================
\section{Introduction}
% ===================================================================

Reliable automated evaluation of AI-generated medical answers is a
prerequisite for safe clinical deployment. The dominant paradigm routes
answers to a single large, closed-source model such as GPT-4~\cite{achiam2023gpt}
acting as an automated judge~\cite{zheng2023judging}. In healthcare,
this approach carries two compounding risks: (i) transmitting protected
health information to third-party APIs violates HIPAA~\cite{cohen2017hipaa},
and (ii) remote inference is unavailable in offline or
bandwidth-constrained clinical environments.

Prior work has explored rubric-grounded judges~\cite{kim2024prometheus},
diverse model panels~\cite{verga2024replacing,chan2023chateval}, and
clinical benchmarking~\cite{singhal2023large,nori2023gpt4}, yet no study
has jointly examined how rubric design, scoring scale, and clinical domain
interact to affect inter-judge agreement in a fully local,
privacy-preserving deployment.

We address this gap with Multi-SLMs-as-Judge, a framework that (1) runs
a configurable panel of locally deployed SLMs, (2) computes a pairwise
agreement score to determine routing, and (3) reports a quality score
only when sufficient consensus is reached. Our main contributions are:
\begin{itemize}
  \item A consensus-gated evaluation framework requiring no external API.
  \item An empirical comparison of five rubrics across three scoring
        scales (binary, Likert~1--5, Likert~0--10) on 100 clinical
        questions in five domains.
  \item Evidence that scoring scale has a substantial independent effect
        on inter-judge agreement, independent of rubric content.
\end{itemize}

% ===================================================================
\section{Related Work}
% ===================================================================

\subsection{LLM-as-Judge}
Zheng et al.\ introduced the LLM-as-a-judge paradigm using MT-Bench
and demonstrated that strong models can serve as reliable evaluators~\cite{zheng2023judging}.
Kim et al.\ showed that fine-grained rubric-grounded prompting substantially
improves consistency~\cite{kim2024prometheus}. Known failure modes include
self-preference bias~\cite{panickssery2024bias} and sensitivity to
prompt framing~\cite{zeng2023framing}.

\subsection{Panel-Based Evaluation}
Verga et al.\ demonstrated that a jury of diverse smaller models can
match single large-judge reliability at lower cost~\cite{verga2024replacing}.
Chan et al.\ used multi-agent debate to improve evaluation
fidelity~\cite{chan2023chateval}. Our framework extends this direction
to the medical domain with an explicit consensus-gated routing mechanism.

\subsection{Medical AI Evaluation}
Singhal et al.\ and Nori et al.\ showed strong benchmark performance
for large models on clinical questions but did not systematically study
evaluation reliability~\cite{singhal2023large,nori2023gpt4}. The
HealthBench rubric, developed with input from 262 physicians, provides
a rigorous framework for safety-critical clinical
assessment~\cite{arora2025healthbench}. Ho et al.\ reviewed qualitative
metrics for evaluating LLMs in clinical decision-making~\cite{ho2024qualitative}.

% ===================================================================
\section{Framework}
% ===================================================================

\subsection{System Overview}

Fig.~\ref{fig:system} shows the Multi-SLMs-as-Judge pipeline.
Each judge $j_i$ in the panel $J = \{J_1, \ldots, J_N\}$ independently
receives three inputs: the medical question $q$, the candidate answer
$a$, and the evaluation rubric $R$. A pairwise agreement score
$\mathrm{Agr}$ is computed first to decide routing. If consensus
exceeds a configurable threshold, a quality score $\hat{S}$ is
reported; otherwise the case is flagged for human review.

\begin{figure}[htbp]
\centerline{\includegraphics[width=\columnwidth]{figures/framework.png}}
\caption{Multi-SLMs-as-Judge pipeline. Each of the $N$ locally deployed
SLMs scores the answer against rubric $R$. The pairwise agreement score
$\mathrm{Agr}$ gates reporting of the quality score $\hat{S}$ and routes
low-consensus cases to human review.}
\label{fig:system}
\end{figure}

\subsection{Pairwise Agreement Score}

Let $s_i^k$ denote the score assigned by judge $i$ to criterion $k$,
and $\max_k$ the maximum achievable score for criterion $k$. The
normalized pairwise agreement score is:

\begin{equation}
\mathrm{Agr} = \frac{2}{N(N-1)}\sum_{i < j}
\left(1 - \frac{1}{|R|}\sum_{k=1}^{|R|}
\frac{|s_i^k - s_j^k|}{\max_k}\right)
\label{eq:agr}
\end{equation}

\noindent where $|R|$ is the number of rubric criteria. Dividing each
criterion difference by $\max_k$ normalizes scores to $[0,1]$, making
binary and Likert rubrics directly comparable. Thus
$\mathrm{Agr} \in [0,1]$, with 1.0 indicating perfect consensus.

Four routing levels are defined by user-configurable thresholds
$\theta_{FA}$, $\theta_{MA}$, and $\theta_{SP}$:

\begin{itemize}
  \item \textbf{Full Agreement (FA):} $\mathrm{Agr} \geq \theta_{FA}$
  \item \textbf{Majority Agreement (MA):} $\theta_{MA} \leq \mathrm{Agr} < \theta_{FA}$
  \item \textbf{Split (SP):} $\theta_{SP} \leq \mathrm{Agr} < \theta_{MA}$
  \item \textbf{Disagreement (D):} $\mathrm{Agr} < \theta_{SP}$
\end{itemize}

Default thresholds used in this study are $\theta_{FA} = 0.95$,
$\theta_{MA} = 0.75$, $\theta_{SP} = 0.50$, calibrated for a four-judge
panel to reflect near-unanimous, supermajority, split, and
substantial-disagreement regimes respectively.

\subsection{Quality Score}

The quality score $\hat{S}$ is reported only when
$\mathrm{Agr} \geq \theta_{MA}$. It is computed as the criterion-weighted
mean of all judges' normalized scores:

\begin{equation}
\hat{S} = \frac{1}{N}\sum_{i=1}^{N}
\frac{\sum_{k=1}^{|R|} w_k \cdot (s_i^k / \max_k)}{\sum_{k=1}^{|R|} w_k}
\label{eq:score}
\end{equation}

\noindent where $w_k$ is the weight of criterion $k$. The resulting
$\hat{S} \in [0,1]$ summarizes how well the answer satisfies the rubric.
When one judge's scores deviate substantially from the rest, the
framework supports three configurable outlier strategies: equal
weighting, outlier removal, or down-weighting.

% ===================================================================
\section{Evaluation Rubrics}
% ===================================================================

We evaluate the framework using five rubrics spanning patient
communication, clinical safety, and general quality assessment.
Table~\ref{tab:rubrics} summarizes the rubric criteria.

\begin{table}[htbp]
\caption{Medical Answer Evaluation Rubrics}
\label{tab:rubrics}
\begin{center}
\begin{tabular}{|l|l|p{4.2cm}|}
\hline
\textbf{ID} & \textbf{Scale} & \textbf{Criteria} \\
\hline
PEMAT~\cite{shoemaker2014development} & Binary & Plain language, organized, focused; actionable steps, barriers \\
\hline
HealthBench~\cite{arora2025healthbench} & Binary & Correctness, safety, escalation, uncertainty, referral \\
\hline
ClinicalEval~\cite{ho2024qualitative} & Likert 1--5 & Accuracy, safety, relevance, completeness, clarity \\
\hline
Prometheus~\cite{kim2024prometheus} & Likert 1--5 & Instruction-following, factuality, coherence, completeness \\
\hline
PEMAT-Likert & Likert 1--5 & Same criteria as PEMAT (scale variant) \\
\hline
\end{tabular}
\end{center}
\end{table}

PEMAT~\cite{shoemaker2014development} assesses patient understandability
and actionability using binary criteria. HealthBench~\cite{arora2025healthbench}
evaluates five safety-critical clinical dimensions. ClinicalEval~\cite{ho2024qualitative}
measures overall clinical utility via a 1--5 Likert scale.
Prometheus~\cite{kim2024prometheus} judges instruction-following,
factual grounding, and coherence. PEMAT-Likert uses identical criteria
to PEMAT but with a 1--5 Likert scale, forming the controlled pair for
isolating the scale effect.

% ===================================================================
\section{Experimental Setup}
% ===================================================================

\subsection{Judge Panel}

The panel consists of four locally deployable SLMs spanning three
distinct base architectures:

\begin{itemize}
  \item MedGemma-4B-IT~\cite{sellergren2025medgemma} (Gemma-3 base)
  \item BioMistral-7B-DARE~\cite{labrak2024biomistral} (Mistral base)
  \item Meditron-7B~\cite{chen2023meditron70b} (LLaMA-2 base)
  \item MedAlpaca-7B~\cite{han2023medalpaca} (LLaMA-2 base)
\end{itemize}

All models were served locally via vLLM with greedy decoding
(temperature~=~0.0, max~1024 output tokens). No data was transmitted
to any external service, satisfying HIPAA data-governance
requirements~\cite{cohen2017hipaa}.

\subsection{Benchmark Dataset}

The evaluation dataset comprises 100 questions drawn from three
publicly available medical QA collections: MedQuAD~\cite{ben-abacha-medquad-2019},
MedDialog~\cite{chen-meddialog-2020}, and Medical
Meadow~\cite{singhal-medical-meadow-2023}. Twenty questions were
sampled from each of five clinical domains: cardiology, pharmacology,
neurology, pediatrics, and emergency medicine. Answers were authored
or verified by medical professionals, reflecting realistic quality
variation including ambiguous and partially correct responses.

\subsection{Hardware}
All experiments were conducted on a high-performance server equipped
with a 64-core CPU and 128~GB of RAM.

% ===================================================================
\section{Results}
% ===================================================================

\subsection{Agreement Across Rubrics (Default Scale)}

Table~\ref{tab:exp1} reports results for all five rubrics at their
default scoring scales. Binary rubrics consistently achieve higher
inter-judge agreement than Likert rubrics. PEMAT and HealthBench
reach 94.1\% and 91.3\% mean agreement, while ClinicalEval, Prometheus,
and PEMAT-Likert fall to 80.9\%, 74.0\%, and 72.2\% respectively.
The direct comparison between PEMAT (binary) and PEMAT-Likert
(identical criteria, Likert scale) isolates the scale effect and
yields a 22 percentage-point agreement drop.

\begin{table}[htbp]
\caption{Results Across All Rubrics at Default Scale}
\label{tab:exp1}
\begin{center}
\small
\begin{tabular}{|l|l|r|r|r|r|r|}
\hline
\textbf{Rubric} & \textbf{Scale} & \textbf{Agr\%} & \textbf{$\hat{S}$} & \textbf{FA} & \textbf{MA} & \textbf{SP+D} \\
\hline
PEMAT        & Binary & $94.1{\pm}10.5$ & $.68{\pm}.12$ & 72 & 24 & 4 \\
HealthBench  & Binary & $91.3{\pm}9.9$  & $.75{\pm}.10$ & 51 & 38 & 11 \\
ClinicalEval & Likert & $80.9{\pm}13.3$ & $.72{\pm}.14$ &  6 & 61 & 33 \\
Prometheus   & Likert & $74.0{\pm}16.2$ & $.55{\pm}.18$ & 17 & 53 & 30 \\
PEMAT-Likert & Likert & $72.2{\pm}8.9$  & $.60{\pm}.11$ &  5 & 50 & 45 \\
\hline
\end{tabular}
\end{center}
\end{table}

\subsection{Effect of Scoring Scale}

Table~\ref{tab:exp2} reports mean pairwise agreement under three
scoring scales: binary, Likert~1--5, and Likert~0--10, applied to each
rubric. Binary elicitation yields the highest agreement for all rubrics
except HealthBench, where the 0--10 scale achieves a marginal gain
(92.6\% vs.\ 91.3\%). Likert~1--5 consistently produces the lowest
agreement. The 0--10 scale partially recovers agreement by allowing
finer score distinctions; ClinicalEval peaks at 96.8\% under binary
scoring, the highest value observed across all experiments.

\begin{table}[htbp]
\caption{Mean Pairwise Agreement (\%) under Three Scoring Scales}
\label{tab:exp2}
\begin{center}
\small
\begin{tabular}{|l|r|r|r|}
\hline
\textbf{Rubric} & \textbf{Binary} & \textbf{Likert 1--5} & \textbf{0--10} \\
\hline
PEMAT        & $\mathbf{94.1}{\pm}10.5$ & $71.5{\pm}10.6$ & $88.6{\pm}4.1$ \\
HealthBench  & $91.3{\pm}9.9$  & $80.8{\pm}14.5$ & $\mathbf{92.6}{\pm}5.7$ \\
ClinicalEval & $\mathbf{96.8}{\pm}9.6$  & $80.9{\pm}13.3$ & $92.1{\pm}5.5$ \\
Prometheus   & $91.3{\pm}17.7$ & $74.0{\pm}16.2$ & $\mathbf{89.4}{\pm}6.4$ \\
PEMAT-Likert & $\mathbf{94.1}{\pm}10.6$ & $72.2{\pm}8.9$  & $88.5{\pm}4.2$ \\
\hline
\end{tabular}
\end{center}
\end{table}

\subsection{Agreement Across Clinical Domains}

Table~\ref{tab:exp3} reports the distribution of agreement levels
across five clinical domains, aggregated over all five rubrics (100
cases per domain). Cardiology is the most consistently judged domain
(36 Full Agreement, 0 Full Disagreement). Pharmacology is the most
contested (20 Full Agreement, 38 combined Split and Full Disagreement),
with elevated disagreement driven primarily by ClinicalEval~(13) and
PEMAT-Likert~(14) on pharmacology questions. PEMAT and HealthBench
show minimal pharmacology contention ($\leq$3 SP+D), confirming that
rubric design strongly moderates the translation of domain difficulty
into panel disagreement.

\begin{table}[htbp]
\caption{Agreement Level Distribution Across Clinical Domains}
\label{tab:exp3}
\begin{center}
\small
\begin{tabular}{|l|r|r|r|r|}
\hline
\textbf{Domain} & \textbf{FA} & \textbf{MA} & \textbf{SP} & \textbf{D} \\
\hline
Cardiology   & 36 & 50 & 14 & 0 \\
Neurology    & 31 & 47 & 22 & 0 \\
Pediatrics   & 31 & 46 & 23 & 0 \\
Emergency    & 33 & 41 & 25 & 1 \\
Pharmacology & 20 & 42 & 33 & 5 \\
\hline
\end{tabular}
\end{center}
\end{table}

% ===================================================================
\section{Discussion}
% ===================================================================

High inter-judge agreement is a necessary but not sufficient condition
for clinical correctness: judges may unanimously assign an incorrect
score when a flaw is subtle. Before any clinical deployment, panels
should be validated against physician ratings on a held-out set.
The default thresholds ($\theta_{FA}=0.95$, $\theta_{MA}=0.75$,
$\theta_{SP}=0.50$) are reasonable for a four-judge panel but should
be recalibrated using domain-specific historical data when panel size
or clinical context changes.

Regarding panel composition, running a rubric-scoring probe and
conducting per-judge ablation experiments before finalizing the panel
are recommended. MedAlpaca's 51.7\% item-leading rate suggests its
agreement contribution may not reflect genuine evaluation, and
full-agreement rates should be treated as upper bounds. Where human
routing costs are prohibitive, a mediator model that reconciles
divergent rationales may reduce escalation rates, though its outputs
require empirical validation against physician adjudication.

% ===================================================================
\section{Conclusion}
% ===================================================================

We presented Multi-SLMs-as-Judge, a privacy-preserving framework for
medical answer evaluation using panels of locally deployed SLMs.
Gating quality score reporting on inter-judge consensus transforms
panel disagreement into an actionable routing signal. Experiments
across five rubrics and five clinical domains show that binary scoring
yields the highest agreement for rubrics with natural yes/no
thresholds, the 0--10 scale most consistently reduces score variance,
and pharmacology generates the most disagreement driven primarily by
Likert-scale rubrics. Because scale and domain effects can be
disentangled, gate thresholds can be tuned independently of rubric
design, making the framework adaptable across clinical e-health
settings. For deployment, we recommend matching scoring scale to rubric
type, probing all candidate judges before panel finalization, and
routing all Split and Disagreement outcomes to human review.

% ===================================================================
\section*{Acknowledgment}
% ===================================================================

The authors thank the developers of MedGemma, BioMistral, Meditron,
and MedAlpaca for making their models publicly available. All
benchmark datasets are used under their respective open licences.

% ===================================================================
% REFERENCES
% ===================================================================
\begin{thebibliography}{00}

\bibitem{achiam2023gpt}
J. Achiam et al., ``GPT-4 technical report,'' arXiv preprint
arXiv:2303.08774, 2023.

\bibitem{arora2025healthbench}
R. K. Arora et al., ``HealthBench: Evaluating large language models
towards improved human health,'' arXiv preprint arXiv:2505.08775, 2025.

\bibitem{ben-abacha-medquad-2019}
A. Ben Abacha and D. Demner-Fushman, ``A question-entailment approach
to question answering,'' in \textit{Proc. 18th BioNLP Workshop}, 2019.

\bibitem{chan2023chateval}
C.-M. Chan et al., ``ChatEval: Towards better LLM-based evaluators
through multi-agent debate,'' arXiv preprint arXiv:2308.07201, 2023.

\bibitem{chen2023meditron70b}
Z. Chen et al., ``MEDITRON-70B: Scaling medical pretraining for large
language models,'' arXiv preprint arXiv:2311.16079, 2023.

\bibitem{chen-meddialog-2020}
S. Chen et al., ``MedDialog: Large-scale medical dialogue datasets,''
arXiv preprint arXiv:2004.03329, 2020.

\bibitem{cohen2017hipaa}
I. G. Cohen and M. M. Mello, ``Unauthorized disclosure of protected
health information: Criminal and civil HIPAA enforcement actions,''
\textit{JAMA}, vol. 318, no. 14, pp. 1339--1340, 2017.

\bibitem{han2023medalpaca}
T. Han et al., ``MedAlpaca: An open-source collection of medical
conversational AI models and training data,'' arXiv preprint
arXiv:2304.08247, 2023.

\bibitem{ho2024qualitative}
C. N. Ho et al., ``Qualitative metrics from the biomedical literature
for evaluating large language models in clinical decision-making:
A narrative review,'' \textit{BMC Medical Informatics and Decision
Making}, vol. 24, no. 1, p. 357, 2024.

\bibitem{kim2024prometheus}
S. Kim et al., ``Prometheus: Inducing fine-grained evaluation capability
in language models,'' in \textit{Proc. ICLR}, 2024.

\bibitem{labrak2024biomistral}
Y. Labrak et al., ``BioMistral: A collection of open-source pretrained
large language models for medical domains,'' arXiv preprint
arXiv:2402.10373, 2024.

\bibitem{nori2023gpt4}
H. Nori et al., ``Capabilities of GPT-4 on medical challenge
problems,'' arXiv preprint arXiv:2303.13375, 2023.

\bibitem{panickssery2024bias}
A. Panickssery, S. R. Bowman, and S. Feng, ``LLM evaluators recognize
and favor their own generations,'' in \textit{Proc. NeurIPS}, vol. 37,
2024.

\bibitem{sellergren2025medgemma}
A. Sellergren et al., ``MedGemma technical report,'' arXiv preprint
arXiv:2507.05201, 2025.

\bibitem{shoemaker2014development}
S. J. Shoemaker, M. S. Wolf, and C. Brach, ``Development of the
Patient Education Materials Assessment Tool (PEMAT): A new instrument
to assess understandability and actionability,'' \textit{Patient
Education and Counseling}, vol. 96, no. 3, pp. 395--403, 2014.

\bibitem{singhal2023large}
K. Singhal et al., ``Large language models encode clinical knowledge,''
\textit{Nature}, vol. 620, pp. 172--180, 2023.

\bibitem{singhal-medical-meadow-2023}
K. Singhal et al., ``Towards expert-level medical question answering
with large language models,'' arXiv preprint arXiv:2305.09617, 2023.

\bibitem{verga2024replacing}
P. Verga et al., ``Replacing judges with juries: Evaluating LLM
generations with a panel of diverse models,'' arXiv preprint
arXiv:2404.18796, 2024.

\bibitem{zeng2023framing}
Z. Zeng et al., ``Evaluating large language models at evaluating
instruction following,'' arXiv preprint arXiv:2310.07641, 2023.

\bibitem{zheng2023judging}
L. Zheng et al., ``Judging LLM-as-a-judge with MT-Bench and chatbot
arena,'' in \textit{Proc. NeurIPS}, vol. 36, 2023.

\end{thebibliography}

\end{document}
